\documentclass[12pt]{article}
\usepackage{amsmath, amssymb, amsthm}

\title{A7 Invariant Layer}
\author{Borealis Serenity Hedling}
\date{28 October 2026}

\begin{document}

\maketitle

\section*{Altitude A7: Invariant Layer}

This module defines the invariant structures governing operator behavior,
collapse-surface geometry, fusion operations, and manifold coupling. These
invariants constrain altitude A8 holonomy-runtime mathematics and ensure
altitude-safe evolution of scalar fields defined in A6.

\section{Operator Invariants}

\begin{definition}[Reversible Operator]
Let $W$ be a world-state space. A \emph{reversible operator} is a map
\[
O : W \to W
\]
such that $O$ is bijective and $O^{-1}$ exists.
\end{definition}

\begin{definition}[Operator Invariant]
An \emph{operator invariant} is a function
\[
I : W \to \mathbb{R}
\]
such that
\[
I(O(W)) = I(W)
\]
for every reversible operator $O$.
\end{definition}

\section{Collapse-Surface Invariants}

\begin{definition}[Collapse-Surface Invariant]
Let $\Sigma$ be a collapse surface associated to a scalar field $\phi$.
A \emph{collapse-surface invariant} is defined by
\[
I_\Sigma(W) = \int_{\Sigma} F(W) \, d\mu,
\]
where $F$ is a smooth observable and $d\mu$ is the induced measure on $\Sigma$.
\end{definition}

\begin{remark}
Collapse-surface invariants encode geometric stability conditions inherited
from altitude A6.
\end{remark}

\section{Fusion Invariants}

\begin{definition}[Fusion Operator]
Given reversible operators $O_1$ and $O_2$, define the \emph{fusion operator}
\[
O_{\mathrm{fuse}} = O_1 \otimes O_2.
\]
\end{definition}

\begin{definition}[Fusion Invariant]
An invariant $I$ is a \emph{fusion invariant} if
\[
I(O_{\mathrm{fuse}}(W)) = I(W)
\]
for all world-states $W$.
\end{definition}

\section{Manifold-Coupling Invariants}

\begin{definition}[Manifold-Coupling Invariant]
Let $T : M \to \mathcal{T}$ be a tensor field on a manifold $M$. A
\emph{manifold-coupling invariant} is defined by
\[
I_M(T) = \int_M \langle T, T \rangle \, d\mu,
\]
where $\langle \cdot, \cdot \rangle$ is a metric-induced inner product.
\end{definition}

\section{Invariant-Preservation Theorems}

\begin{theorem}[Operator Invariant Preservation]
Let $O$ be a reversible operator. Then
\[
I(O(W)) = I(W).
\]
\end{theorem}

\begin{proof}[Proof Sketch]
Follows from the definition of operator invariants. Full formalization appears
in the Lean module \texttt{PreservationProofs.lean}.
\end{proof}

\begin{theorem}[Collapse-Surface Compatibility]
Let $\Sigma$ be a collapse surface associated to $\phi$. Then
\[
I_\Sigma(O(W)) = I_\Sigma(W)
\]
for all reversible operators $O$ that preserve $\Sigma$.
\end{theorem}

\section{Altitude-Safety Constraints}

\begin{definition}[Altitude-Safety Constraints]
Invariant structures must satisfy:
\begin{enumerate}
    \item Operator invariants must be preserved under all reversible operators.
    \item Collapse-surface invariants must respect scalar-field geometry from A6.
    \item Fusion invariants must remain stable under operator tensor products.
    \item Manifold-coupling invariants must be metric-consistent.
\end{enumerate}
\end{definition}

\section{Future Formalization}

\begin{itemize}
    \item Lean formalization of operator invariants.
    \item Lean proofs of collapse-surface invariants.
    \item Lean encoding of fusion invariants.
    \item Integration with A8 holonomy constraints.
\end{itemize}

\end{document}
